\documentclass[10pt]{article}
\usepackage[T1]{fontenc}
\usepackage[utf8]{inputenc}
\usepackage[french]{babel}
\usepackage{amsmath,amssymb,mathrsfs,mathpazo}
\usepackage{geometry}
\geometry{a4paper,margin=2cm}
\usepackage{graphicx,array,tabularx,tabularray,colortbl,multirow,multicol,booktabs,enumitem,pifont,ulem,fancybox,diagbox,cancel}
\usepackage{pstricks,pstricks-add,pst-plot,pst-tree,pst-node,pst-text,pst-func,pst-math,pst-bspline,pst-3dplot,pst-eucl,pst-solides3d}
\let\pscalculate\undefined
\usepackage{pst-bezier}
\usepackage{tikz,pgfplots,tkz-tab}
\usetikzlibrary{arrows,arrows.meta,patterns,calc,positioning,trees,decorations.markings,shapes.geometric,shapes.misc,shapes.arrows,intersections,decorations.pathmorphing,decorations.pathreplacing,backgrounds,shadows,babel}
\pgfplotsset{compat=1.18}
\usepgfplotslibrary{fillbetween}
\usepackage[np]{numprint}
\usepackage{listings,makecell,pas-tableur,twoopt}
\let\barre\undefined
\usepackage[pstricks]{tablvar}
\definecolor{darkgreen}{RGB}{0,100,0}
\providecommand{\red}{\color{red}}
\providecommand{\blue}{\color{blue}}
\providecommand{\green}{\color{green}}
\providecommand{\black}{\color{black}}
\providecommand{\white}{\color{white}}
\providecommand{\gray}{\color{gray}}
\providecommand{\cyan}{\color{cyan}}
\providecommand{\magenta}{\color{magenta}}
\providecommand{\yellow}{\color{yellow}}
\providecommand{\violet}{\color{violet}}
\providecommand{\orange}{\color{orange}}
\providecommand{\euro}{\texteuro}
\providecommand{\EUR}{\texteuro}
\providecommand{\ieme}{\textsuperscript{e}}
\providecommand{\degres}{\ensuremath{^\circ}}
\providecommand{\degree}{\ensuremath{^\circ}}
\providecommand{\vv}[1]{\overrightarrow{#1}}
\providecommand{\Coord}[1]{#1}
\providecommand{\decosix}{}
\providecommand{\starredbullet}{\ensuremath{\bullet}}
\pagestyle{empty}
\setlength{\parindent}{0pt}
\usepackage{adjustbox}\DeclareUnicodeCharacter{2236}{:}\DeclareUnicodeCharacter{2190}{\ensuremath{\leftarrow}}
\providecommand{\R}{\mathbb{R}}
\providecommand{\N}{\mathbb{N}}
\providecommand{\Z}{\mathbb{Z}}
\providecommand{\Q}{\mathbb{Q}}
\providecommand{\D}{\mathbb{D}}
\providecommand{\C}{\mathbb{C}}
\providecommand{\e}{{\mathrm{e}}}
\providecommand{\vect}[1]{\overrightarrow{#1}}
\providecommand{\vectt}[1]{\overrightarrow{\mathrm{#1}}}
\providecommand{\vvt}[1]{\overrightarrow{\mathrm{#1}}}
\providecommand{\np}[2][]{{#2~#1}}
\providecommand{\vv}[1]{{\overrightarrow{#1}}}
\renewcommand{\d}{{\mathrm{d}}}
\providecommand{\degres}{{^\circ}}
\providecommand{\degre}{{^\circ}}
\providecommand{\blue}{\color{blue}}
\providecommand{\red}{\color{red}}
\providecommand{\enskip}{\quad}
\providecommand{\llbracket}{\lbrack\!\lbrack}
\providecommand{\rrbracket}{\rbrack\!\rbrack}
\providecommand{\euro}{{\text{€}}}
\providecommand{\no}{{\text{n°}}}
\providecommand{\interval}[2]{{[#1;#2]}}
\providecommand{\intervFO}[2]{{[#1;#2[}}
\providecommand{\dotfill}{\ldots}

\providecommand{\ds}{\displaystyle}
\providecommand{\textding}[1]{#1}
\providecommand{\Oij}{\(\mathrm{O};\vec{\imath},\vec{\jmath}\)}
\providecommand{\Oijk}{\(\mathrm{O};\vec{\imath},\vec{\jmath},\vec{k}\)}
\providecommand{\Ouv}{\(\mathrm{O};\vec{u},\vec{v}\)}
\providecommand{\intervFF}[2]{\ensuremath{[#1;#2]}}
\renewcommand{\euro}{\ifmmode\text{\texteuro}\else\texteuro\fi}
\DeclareUnicodeCharacter{00B0}{\ifmmode^{\circ}\else\textdegree\fi}
\DeclareUnicodeCharacter{20AC}{\ifmmode\text{\texteuro}\else\texteuro\fi}
% Énoncés : Éducation nationale. Transcription et corrigés : APMEP, Denis Vergès et contributeurs. Découpage et classement : CoopMaths, Sébastien Lozano, Éric Elter et contributeurs. Adaptation de la présentation : Maths Explorer.
% CC BY-SA : https://coopmaths.fr/www/about/
% https://coopmaths.fr/alea/static/bac/2021/tex/bac_2021_06_asie_5.tex
\begin{document}
\noindent\textit{APMEP / CoopMaths. CC BY-SA. Présentation adaptée par Maths Explorer.}\par\medskip
\medskip

\begin{tabular}{|l|}\hline
\textbf{Principaux domaines abordés}\\
-- Étude de fonction, fonction exponentielle\\
-- Équations différentielles\\ \hline
\end{tabular}

\bigskip

\textbf{Partie I}

\medskip

Considérons l'équation différentielle

\[y'= -0,4y + 0,4\]

où $y$ désigne une fonction de la variable $t$, définie et dérivable sur $[0~;~ + \infty[$.

\medskip

\begin{enumerate}
\item 
	\begin{enumerate}
		\item Déterminer une solution particulière constante de cette équation différentielle. 
		\item En déduire l'ensemble des solutions de cette équation différentielle.
\item Déterminer la fonction $g$, solution de cette équation différentielle, qui vérifie $g(0) = 10$.
	\end{enumerate}
\end{enumerate}

\bigskip

\textbf{Partie II}

\medskip

Soit $p$ la fonction définie et dérivable sur l'intervalle $[0~;~+ \infty[$ par 

\[p(t) = \dfrac{1}{g(t)} = \dfrac{1}{1 + 9\e^{-0,4t}}.\]

\smallskip

\begin{enumerate}
\item Déterminer la limite de $p$ en $+ \infty$. 
\item Montrer que $p'(t) = \dfrac{3,6\e^{-0,4t}}{ \left(1 + 9\e^{-0,4t}\right)^2}$ pour tout $t \in  [0~;~+ \infty[$.

\item  
	\begin{enumerate}
		\item Montrer que l'équation $p(t) = \dfrac{1}{2}$ admet une unique solution $\alpha$ sur $[0~;~+ \infty[$. 
		\item Déterminer une valeur approchée de $\alpha$ à $10^{-1}$ près à l'aide d'une calculatrice.
	\end{enumerate}
\end{enumerate}

\bigskip

\textbf{Partie III}

\medskip

\begin{enumerate}
\item $p$ désigne la fonction de la partie II.

Vérifier que $p$ est solution de l'équation différentielle $y' = 0,4y(1 - y)$ avec la condition initiale 
$y(0) = \dfrac{1}{10}$ où $y$ désigne une fonction définie et dérivable sur $[0~;~ + \infty[$.
\item Dans un pays en voie de développement, en l'année 2020, 10\,\% des écoles ont accès à internet. 

Une politique volontariste d'équipement est mise en œuvre et on s'intéresse à l'évolution de la proportion des écoles ayant accès à internet. 

On note $t$ le temps écoulé, exprimé en année, depuis l'année 2020.

La proportion des écoles ayant accès à internet à l'instant $t$ est modélisée par $p(t)$.

Interpréter dans ce contexte la limite de la question II 1 puis la valeur approchée de $\alpha$ de la question II 3. b. ainsi que la valeur $p(0)$.
\end{enumerate}

\end{document}
